Pour mesurer la longueur d'onde d'un laser, une étudiante éclaire deux fentes séparées de \qty{\v{sep}}{mm} et observe la figure d'interférence sur un mur situé à \qty{\v{L}}{m}. Elle mesure une distance de \qty{\v{dy}}{mm} entre le centre de la frange brillante centrale et le centre de la frange brillante d'ordre \v{N}.

\begin{enumerate}[label=\alph*)]
    \item Déterminez la longueur d'onde du laser.
    \item À quelle distance du centre se trouve la deuxième frange sombre ?
    \item On plonge tout le montage dans l'eau ($n = \num{1.33}$). Quelle est la nouvelle distance entre deux franges brillantes voisines ?
\end{enumerate}
